\begin{proof}[Proof of \cref{obj:lem:jackson-tensor-extraction}]
\begin{enumerate}
\item Define
\[
\operatorname{Box}_x
=
\{y\in\mathbb R^{\mathcal J}: |y_\jmath-b_{\jmath,x}|\le r_{\jmath,x}
\ \text{for every }\jmath\in\mathcal J\},
\qquad
\operatorname{Cube}_{\mathcal J}=[-1,1]^{\mathcal J},
\]
and
\[
\operatorname{aff}_x(\zeta)_\jmath
=
b_{\jmath,x}+r_{\jmath,x}\zeta_\jmath,
\qquad
\zeta\in\mathbb R^{\mathcal J}.
\]
If \(y\in\operatorname{Box}_x\), then
\[
y_\jmath\ge b_{\jmath,x}-r_{\jmath,x}
=
\frac{\ell_{\jmath,x}+u_{\jmath,x}}2
-
\frac{u_{\jmath,x}-\ell_{\jmath,x}}2
=
\ell_{\jmath,x}\ge0
\qquad(\jmath\in\mathcal J),
\]
and \(b_x\in\operatorname{Box}_x\), since each \(r_{\jmath,x}>0\).

For nonnegative four-coordinate vectors \(u,w\), put
\[
\delta_{\mathcal J}(u,w)
=
\sum_{\jmath\in\mathcal J}|u_\jmath-w_\jmath|.
\]
For each arm \(a\in\{0,1\}\), \cref{obj:def:global-extension} gives
\[
g_{a,\epsilon}(u)
=
\begin{cases}
0, & s(u)=0,\\[2mm]
\dfrac{s(u)u_{a1}}{\max\{s_a(u),\epsilon s(u)\}}, & s(u)>0.
\end{cases}
\]
Moreover
\[
0\le g_{a,\epsilon}(u)\le s(u),
\]
because \(u_{a1}\le s_a(u)\le \max\{s_a(u),\epsilon s(u)\}\) when \(s(u)>0\), and the zero-mass value is \(0\). If \(s(u)=0\), nonnegativity gives \(u=0\), hence
\[
|g_{a,\epsilon}(u)-g_{a,\epsilon}(w)|
=
g_{a,\epsilon}(w)
\le s(w)
=
\delta_{\mathcal J}(u,w)
\le (1+\epsilon^{-1})\delta_{\mathcal J}(u,w).
\]
The case \(s(w)=0\) is symmetric.

Assume now that \(s(u)>0\) and \(s(w)>0\). The elementary bounds
\[
|s(u)-s(w)|\le\delta_{\mathcal J}(u,w),
\qquad
|u_{a1}-w_{a1}|\le\delta_{\mathcal J}(u,w),
\]
and
\[
|s_a(u)-s_a(w)|
\le
|u_{a0}-w_{a0}|+|u_{a1}-w_{a1}|
\le
\delta_{\mathcal J}(u,w)
\]
will be used throughout. If
\[
s_a(u)\le\epsilon s(u),
\qquad
s_a(w)\le\epsilon s(w),
\]
then
\[
g_{a,\epsilon}(u)=\frac{u_{a1}}{\epsilon},
\qquad
g_{a,\epsilon}(w)=\frac{w_{a1}}{\epsilon},
\]
so
\[
|g_{a,\epsilon}(u)-g_{a,\epsilon}(w)|
\le
\epsilon^{-1}\delta_{\mathcal J}(u,w)
\le
(1+\epsilon^{-1})\delta_{\mathcal J}(u,w).
\]
If
\[
\epsilon s(u)\le s_a(u),
\qquad
\epsilon s(w)\le s_a(w),
\]
then
\[
g_{a,\epsilon}(u)=s(u)\frac{u_{a1}}{s_a(u)},
\qquad
g_{a,\epsilon}(w)=s(w)\frac{w_{a1}}{s_a(w)}.
\]
Since \(0\le w_{a1}/s_a(w)\le1\) and \(s(u)/s_a(u)\le\epsilon^{-1}\), while
\[
\left|
\frac{u_{a1}}{s_a(u)}
-
\frac{w_{a1}}{s_a(w)}
\right|
\le
\frac{|u_{a1}-w_{a1}|+|u_{a0}-w_{a0}|}{s_a(u)},
\]
we obtain
\[
\begin{aligned}
|g_{a,\epsilon}(u)-g_{a,\epsilon}(w)|
&\le
|s(u)-s(w)|
+
s(u)
\left|
\frac{u_{a1}}{s_a(u)}
-
\frac{w_{a1}}{s_a(w)}
\right|\\
&\le
\delta_{\mathcal J}(u,w)
+
\epsilon^{-1}\delta_{\mathcal J}(u,w)
=
(1+\epsilon^{-1})\delta_{\mathcal J}(u,w).
\end{aligned}
\]
If
\[
s_a(u)\le\epsilon s(u),
\qquad
\epsilon s(w)\le s_a(w),
\]
then
\[
g_{a,\epsilon}(u)-g_{a,\epsilon}(w)
=
\frac{u_{a1}-w_{a1}}{\epsilon}
+
\frac{w_{a1}}{s_a(w)}
\frac{s_a(w)-\epsilon s(w)}{\epsilon}.
\]
The second summand is nonnegative, giving
\[
g_{a,\epsilon}(u)-g_{a,\epsilon}(w)
\ge
-\epsilon^{-1}\delta_{\mathcal J}(u,w).
\]
For the upper bound, \(0\le w_{a1}/s_a(w)\le1\) and
\[
0\le s_a(w)-\epsilon s(w)
\le
\{s_a(w)-s_a(u)\}-\epsilon\{s(w)-s(u)\},
\]
so
\[
\begin{aligned}
g_{a,\epsilon}(u)-g_{a,\epsilon}(w)
&\le
\frac{u_{a1}-w_{a1}+s_a(w)-s_a(u)-\epsilon\{s(w)-s(u)\}}{\epsilon}\\
&=
\epsilon^{-1}(w_{a0}-u_{a0})-\{s(w)-s(u)\}\\
&\le
(1+\epsilon^{-1})\delta_{\mathcal J}(u,w).
\end{aligned}
\]
The remaining mixed case,
\[
\epsilon s(u)\le s_a(u),
\qquad
s_a(w)\le\epsilon s(w),
\]
is analogous. Here
\[
g_{a,\epsilon}(u)-g_{a,\epsilon}(w)
=
\frac{u_{a1}-w_{a1}}{\epsilon}
-
\frac{u_{a1}}{s_a(u)}
\frac{s_a(u)-\epsilon s(u)}{\epsilon}.
\]
The subtracted summand is nonnegative, hence
\[
g_{a,\epsilon}(u)-g_{a,\epsilon}(w)
\le
\epsilon^{-1}\delta_{\mathcal J}(u,w).
\]
Also \(0\le u_{a1}/s_a(u)\le1\) and
\[
0\le s_a(u)-\epsilon s(u)
\le
\{s_a(u)-s_a(w)\}-\epsilon\{s(u)-s(w)\},
\]
which yields
\[
\begin{aligned}
g_{a,\epsilon}(u)-g_{a,\epsilon}(w)
&\ge
\frac{u_{a1}-w_{a1}-s_a(u)+s_a(w)+\epsilon\{s(u)-s(w)\}}{\epsilon}\\
&=
\epsilon^{-1}(w_{a0}-u_{a0})+\{s(u)-s(w)\}\\
&\ge
-(1+\epsilon^{-1})\delta_{\mathcal J}(u,w).
\end{aligned}
\]
Thus, for both arms,
\[
|g_{a,\epsilon}(u)-g_{a,\epsilon}(w)|
\le
(1+\epsilon^{-1})\delta_{\mathcal J}(u,w).
\]
Since
\[
|\max\{\rho_0,\rho_1\}-\max\{\sigma_0,\sigma_1\}|
\le
\max_{a\in\{0,1\}}|\rho_a-\sigma_a|,
\]
\cref{obj:def:global-extension} gives
\[
|\bar f_\epsilon(u)-\bar f_\epsilon(w)|
\le
(1+\epsilon^{-1})\delta_{\mathcal J}(u,w).
\]
Consequently \(\bar f_\epsilon\) is continuous on \(\operatorname{Box}_x\), and for every
\(y\in\operatorname{Box}_x\),
\[
|\bar f_\epsilon(y)-\bar f_\epsilon(b_x)|
\le
(1+\epsilon^{-1})\sum_{\jmath\in\mathcal J}|y_\jmath-b_{\jmath,x}|
\le
(1+\epsilon^{-1})\sum_{\jmath\in\mathcal J}r_{\jmath,x}.
\]
Set
\[
S_x
=
(1+\epsilon^{-1})\sum_{\jmath\in\mathcal J}r_{\jmath,x}.
\]
% lean: globalCellValue_lipschitz

\item Define the centered pullback on the normalized cube by
\[
\chi_x(\zeta)
=
\bar f_\epsilon(\operatorname{aff}_x(\zeta))-\bar f_\epsilon(b_x),
\qquad
\zeta\in\operatorname{Cube}_{\mathcal J}.
\]
The affine map sends \(\operatorname{Cube}_{\mathcal J}\) into \(\operatorname{Box}_x\), so \(\chi_x\) is continuous and
\[
|\chi_x(\zeta)|\le S_x
\qquad(\zeta\in\operatorname{Cube}_{\mathcal J}).
\]
Since \(K\ge2\), we have \(K\ge1\). Applying \cref{obj:lem:tensor-jackson-coefficient-envelope} to \(\chi_x\) with envelope \(B=S_x\) gives a polynomial \(\mathcal R_{x,K}\) in the normalized coordinates such that
\[
\mathcal R_{x,K}(\cos\vartheta)
=
\int_{[-\pi,\pi]^{\mathcal J}}
\chi_x\{\cos(\vartheta-\tau)\}
\prod_{\jmath\in\mathcal J}J_K(\tau_\jmath)\,d\tau
\qquad(\vartheta:\mathcal J\to\mathbb R),
\]
\[
\beta_\omega\le 2(K-1)
\quad
\text{for every }\beta\in\operatorname{supp}(\mathcal R_{x,K})
\text{ and }\omega\in\mathcal J,
\]
and
\[
\sum_{\beta\in\operatorname{supp}(\mathcal R_{x,K})}
|(\mathcal R_{x,K})_\beta|
\le
2^{40K+20}S_x .
\]
% lean: tensorConvolution_exists_mvPolynomial_four_coeffBound

\item Substitute normalized coordinates into physical coordinates. Since every \(r_{\jmath,x}>0\), there is a polynomial \(P^0_{x,K}\) in the four physical coordinates satisfying
\[
P^0_{x,K}(v)
=
\mathcal R_{x,K}
\left(
\left\{\frac{v_\jmath-b_{\jmath,x}}{r_{\jmath,x}}\right\}_{\jmath\in\mathcal J}
\right)
\qquad(v\in\mathbb R^{\mathcal J}),
\]
and every coordinate degree of \(P^0_{x,K}\) is at most \(2(K-1)\). Define
\[
P_{x,K}(v)=P^0_{x,K}(v)+\bar f_\epsilon(b_x).
\]
Adding a constant preserves the coordinate-degree bound, and therefore
\[
P_{x,K}(v)-\bar f_\epsilon(b_x)
=
\mathcal R_{x,K}
\left(
\left\{\frac{v_\jmath-b_{\jmath,x}}{r_{\jmath,x}}\right\}_{\jmath\in\mathcal J}
\right)
\qquad(v\in\mathbb R^{\mathcal J}).
\]
% lean: mvPolynomial_affine_substitution_four

\item For an angle vector \(\vartheta:\mathcal J\to\mathbb R\),
\[
\left\{
\frac{b_{\jmath,x}+r_{\jmath,x}\cos(\vartheta_\jmath)-b_{\jmath,x}}
{r_{\jmath,x}}
\right\}_{\jmath\in\mathcal J}
=
\{\cos(\vartheta_\jmath)\}_{\jmath\in\mathcal J}.
\]
Thus
\[
P_{x,K}(b_x+r_x\odot\cos\vartheta)
=
\mathcal R_{x,K}(\cos\vartheta)+\bar f_\epsilon(b_x).
\]
Using the convolution identity for \(\mathcal R_{x,K}\),
\[
\begin{aligned}
P_{x,K}(b_x+r_x\odot\cos\vartheta)
&=
\int_{[-\pi,\pi]^{\mathcal J}}
\{\bar f_\epsilon(b_x+r_x\odot\cos(\vartheta-\tau))-\bar f_\epsilon(b_x)\}
\prod_{\jmath\in\mathcal J}J_K(\tau_\jmath)\,d\tau\\
&\qquad+\bar f_\epsilon(b_x).
\end{aligned}
\]
The integrands are continuous on the compact period box, and the tensor Jackson kernel has unit mass,
\[
\int_{[-\pi,\pi]^{\mathcal J}}\prod_{\jmath\in\mathcal J}J_K(\tau_\jmath)\,d\tau=1.
\]
Therefore the constant centered term integrates exactly to the value added back:
\[
\begin{aligned}
&\int_{[-\pi,\pi]^{\mathcal J}}
\{\bar f_\epsilon(b_x+r_x\odot\cos(\vartheta-\tau))-\bar f_\epsilon(b_x)\}
\prod_{\jmath\in\mathcal J}J_K(\tau_\jmath)\,d\tau
+
\bar f_\epsilon(b_x) \\
&\qquad=
\int_{[-\pi,\pi]^{\mathcal J}}
\bar f_\epsilon(b_x+r_x\odot\cos(\vartheta-\tau))
\prod_{\jmath\in\mathcal J}J_K(\tau_\jmath)\,d\tau
-
\bar f_\epsilon(b_x)
\int_{[-\pi,\pi]^{\mathcal J}}\prod_{\jmath\in\mathcal J}J_K(\tau_\jmath)\,d\tau
+
\bar f_\epsilon(b_x) \\
&\qquad=
\int_{[-\pi,\pi]^{\mathcal J}}
\bar f_\epsilon(b_x+r_x\odot\cos(\vartheta-\tau))
\prod_{\jmath\in\mathcal J}J_K(\tau_\jmath)\,d\tau .
\end{aligned}
\]
Thus
\[
P_{x,K}(b_x+r_x\odot\cos\vartheta)
=
\int_{[-\pi,\pi]^{\mathcal J}}
\bar f_\epsilon(b_x+r_x\odot\cos(\vartheta-\tau))
\prod_{\jmath\in\mathcal J}J_K(\tau_\jmath)\,d\tau.
\]
For each coordinate, \(J_K(-t)=J_K(t)\): both sine factors in its defining ratio change sign, and the ratio is raised to the fourth power. The coordinatewise substitution \(t=-\tau\) preserves \([-\pi,\pi]^{\mathcal J}\) and Lebesgue measure, and therefore changes \(\cos(\vartheta-\tau)\) to \(\cos(\vartheta+t)\) while leaving the product kernel unchanged. The last display is consequently the order-\(K\) tensor Jackson convolution in \cref{obj:def:jackson-kernel} of \(\bar f_\epsilon\) precomposed with the affine parametrization of \(Q_x\), evaluated at \(\vartheta\).
% lean: tensorJackson_integral_eq_one

\item Substituting the definition of \(S_x\) into the coefficient estimate gives
\[
\sum_{\beta\in\operatorname{supp}(\mathcal R_{x,K})}
|(\mathcal R_{x,K})_\beta|
\le
2^{40K+20}
(1+\epsilon^{-1})
\sum_{\jmath\in\mathcal J}r_{\jmath,x}.
\]
The constructed \(P_{x,K}\) and \(\mathcal R_{x,K}\) have the asserted physical degree, normalized degree, centered-evaluation identity, Jackson-convolution identity, and coefficient envelope.
% lean: jacksonTensorPolynomialData_exists
\end{enumerate}
\end{proof}
